\section{Análisis de problemas: patrones de falla típicos de los Agentes}

Antes de profundizar en las soluciones, es necesario comprender primero cómo fallan específicamente los Agentes en escenarios de ejecución prolongada. El equipo de Anthropic observó durante sus experimentos dos patrones de falla principales.

\subsection{Patrón de falla uno: intentar completar todas las tareas de una sola vez}

Incluso al utilizar la tecnología de compaction (compresión del contexto), los modelos de vanguardia como Opus 4.5 tienden aún a intentar completar toda la aplicación de un solo golpe frente a prompts de alto nivel (como "construir un clon de claude.ai"). Este comportamiento "one-shot" conduce a:

\begin{itemize}
    \item El Agente agota el contexto (context window) durante la implementación.
    \item Solo se completan las funciones a la mitad y no queda documentación registrada.
    \item El Agente del siguiente sesión necesita gastar mucho tiempo adivinando qué sucedió antes.
    \item El Agente intenta reiteradamente ejecutar de nuevo las funciones básicas, desperdiciando muchos tokens.
\end{itemize}

\begin{warningbox}{La compaction no es una panacea}
La compaction es la capacidad de gestión del contexto del SDK para Agentes de Claude; puede comprimir información histórica cuando el contexto (context window) está cerca de su límite máximo. Sin embargo, la compaction no siempre logra transmitir instrucciones totalmente claras al siguiente Agente: la información comprimida podría perder detalles cruciales, lo que impide que los Agentes posteriores restauren con precisión el estado de trabajo previo. No se debe depender únicamente de la compaction para resolver problemas de continuidad entre sesiones.
\end{warningbox}

\subsection{Patrón de falla dos: anunciar la finalización de la tarea prematuramente}

El segundo patrón de falla suele aparecer en las etapas finales del proyecto. Cuando algunas funciones ya han sido construidas, las instancias posteriores del Agente revisan el estado actual del proyecto; al observar que ya existe cierto progreso, declaran directamente que toda la tarea está completada ("declare the job done"), cuando en realidad quedan muchas funciones sin implementar.

\subsection{Descomposición de la esencia del problema}

Estos dos patrones de falla pueden descomponerse en dos subproblemas:

\begin{table}[H]
\centering
\begin{tabular}{p{4cm}p{9cm}}
\toprule
\textbf{Subproblema} & \textbf{Descripción} \\
\midrule
Inicialización del entorno & Se necesita construir la arquitectura base para todas las funciones pendientes de implementar, guiando al Agente para que avance paso a paso y por función \\
\midrule
Avance incremental & Cada sesión del Agente debería lograr avances incrementales y dejar el entorno en un "estado limpio" al finalizar \\
\bottomrule
\end{tabular}
\caption{Dos subproblemas centrales que enfrentan los Agentes de ejecución prolongada}
\end{table}

\begin{importantbox}{¿Qué es un "estado limpio"?}
Un "estado limpio" (clean state) se refiere a que el código ha alcanzado el estándar para integrarse en la rama principal: no tiene errores graves, está bien organizado y cuenta con buena documentación; los nuevos desarrolladores (o las nuevas sesiones del Agente) pueden comenzar inmediatamente el desarrollo de la siguiente función sin necesidad de limpiar primero el desorden dejado por trabajos anteriores. Esto es totalmente consistente con las mejores prácticas de los ingenieros humanos.
\end{importantbox}

\subsection{Resumen de la sección}

Los Agentes presentan dos patrones de falla típicos en escenarios de ejecución prolongada: intentar devorar más de lo que pueden al querer completar todas las tareas de una sola vez, y anunciar prematuramente la finalización del proyecto tras cierto avance. La raíz de ambos problemas radica en la falta de mecanismos estructurados para la descomposición de tareas y el seguimiento del progreso.

